\documentclass[11pt,a4paper]{article}
\usepackage{fontspec}
\usepackage{polyglossia}
\setmainlanguage{russian}
\setotherlanguage{english}
\setmainfont{Liberation Serif}
\usepackage[left=1.0cm,right=1.0cm,top=0.8cm,bottom=0.7cm]{geometry}
\usepackage{amsmath,amssymb}
\usepackage{array}
\usepackage{multirow}
\setlength{\parindent}{0pt}
\pagestyle{empty}

\newcounter{zad}

% ФОРМАТ. Поле ответа слева, условие справа; верх рамки поля стоит на одной
% горизонтальной линии с верхом строки условия (подгонка — \podnyat).
% Полоска под плюс ОДНА на задачу: плюс ставится за задачу целиком.
% Значок уровня только у обязательных: ○ круг, △ треугольник.
% Дополнительные идут второй страницей и значков не несут.
% Вариант Б — изоморфизм A: та же абстрактная форма, другие числа и сюжеты.
% Ответы обоих вариантов считает otvety.py; все они попарно различны.

\newlength{\otvsh}\setlength{\otvsh}{3.5cm}
\newlength{\plsh}\setlength{\plsh}{0.7cm}
\newlength{\polesh}\setlength{\polesh}{4.6cm}
\newlength{\uslsh}\setlength{\uslsh}{13.6cm}
\newlength{\podnyat}\setlength{\podnyat}{1.1ex}
\newlength{\minpole}\setlength{\minpole}{1.3cm}

\newsavebox{\uslbox}
\newlength{\uslvys}
\newlength{\blh}

\newcommand{\calcblh}[1]{%
  \setlength{\blh}{\dimexpr(\uslvys-\arrayrulewidth)/#1-\arrayrulewidth\relax}}
\newcommand{\mt}[1]{\parbox[t][\blh][t]{\otvsh}{\small\bfseries #1\strut}}

\newcommand{\poleA}{\calcblh{1}\setlength{\tabcolsep}{3pt}%
  \begin{tabular}{|p{\otvsh}|p{\plsh}|}\hline \mt{} & \\ \hline\end{tabular}}
\newcommand{\poleB}[2]{\calcblh{2}\setlength{\tabcolsep}{3pt}%
  \begin{tabular}{|p{\otvsh}|p{\plsh}|}\hline
  \mt{#1} & \multirow{2}{*}{} \\ \cline{1-1} \mt{#2} & \\ \hline\end{tabular}}
\newcommand{\poleC}[3]{\calcblh{3}\setlength{\tabcolsep}{3pt}%
  \begin{tabular}{|p{\otvsh}|p{\plsh}|}\hline
  \mt{#1} & \multirow{3}{*}{} \\ \cline{1-1} \mt{#2} & \\ \cline{1-1}
  \mt{#3} & \\ \hline\end{tabular}}

\newcommand{\poleD}[4]{\calcblh{4}\setlength{\tabcolsep}{3pt}%
  \begin{tabular}{|p{\otvsh}|p{\plsh}|}\hline
  \mt{#1} & \multirow{4}{*}{} \\ \cline{1-1} \mt{#2} & \\ \cline{1-1}
  \mt{#3} & \\ \cline{1-1} \mt{#4} & \\ \hline\end{tabular}}

\newcommand{\poleshirW}[1]{\setlength{\blh}{#1}\setlength{\tabcolsep}{3pt}%
  \begin{tabular}{|p{\dimexpr\uslsh+\polesh-2\tabcolsep-2\arrayrulewidth\relax}|}
  \hline \parbox[t][\blh][t]{\dimexpr\uslsh+\polesh-2\tabcolsep-2\arrayrulewidth\relax}{\strut} \\ \hline\end{tabular}}

\newlength{\mezhzad}\setlength{\mezhzad}{14pt}
\newcommand{\znak}[1]{\ifx&#1&\else$#1$\ \fi}

\newcommand{\zad}[3]{%
  \stepcounter{zad}%
  \sbox{\uslbox}{\parbox[t]{\uslsh}{\strut\textbf{\znak{#1}\thezad.}\ #3\strut}}%
  \setlength{\uslvys}{\dimexpr\ht\uslbox+\dp\uslbox\relax}%
  \ifdim\uslvys<\minpole\setlength{\uslvys}{\minpole}\fi
  \par\vspace{\mezhzad}\noindent
  \begin{minipage}[t]{\polesh}\vspace{0pt}\vspace*{-\podnyat}#2\end{minipage}%
  \hfill
  \begin{minipage}[t]{\uslsh}\vspace{0pt}\usebox{\uslbox}\end{minipage}\par}

\newcommand{\zadw}[3]{%
  \stepcounter{zad}%
  \par\vspace{\mezhzad}\noindent
  \parbox{\dimexpr\uslsh+\polesh\relax}{\strut\textbf{\znak{#1}\thezad.}\ #3\strut}%
  \par\vspace{3pt}\noindent\poleshirW{#2}\par}

\newcommand{\zadd}[2]{%
  \stepcounter{zad}%
  \par\vspace{\mezhzad}\noindent
  \parbox{\dimexpr\uslsh+\polesh\relax}{\strut\textbf{\znak{#1}\thezad.}\ #2\strut}\par}

\newcommand{\razdelt}[1]{\par\noindent{\bfseries\LARGE #1}\par\vspace{2pt}}

\newcommand{\shapka}[1]{%
  \setlength{\mezhzad}{14pt}\setlength{\minpole}{1.3cm}%
  \noindent{\LARGE\bfseries Комбинаторика~1}%
  \hfill{\small Спецмат, 7\,И\quad·\quad22 сентября\quad·\quad вариант~#1}%
  \vspace{4pt}\hrule\vspace{10pt}
  \noindent Фамилия, имя:\ \hrulefill\rule[-3mm]{0pt}{8mm}\par
  \vspace{13pt}
  \noindent{\bfseries\LARGE Обязательные задачи}\par\vspace{2pt}
  \setcounter{zad}{0}}

\newcommand{\vtoraya}{\newpage
  \setlength{\mezhzad}{8pt}\setlength{\minpole}{0.85cm}
  \razdelt{Дополнительные задачи}}

\begin{document}

%%%%%%%%%%%%%%%%%%%%%%%%%%%% ВАРИАНТ A %%%%%%%%%%%%%%%%%%%%%%%%%%%%
\shapka{A}

\zad{\bigcirc}{\poleC{а}{б}{в}}{\textbf{(а)}~Сколько чисел в ряду $17,\ 18,\ 19,\ \ldots,\ 122,\ 123$?\newline
\textbf{(б)}~Сколько чисел в ряду $-17,\ -16,\ \ldots,\ -1,\ 0,\ 1,\ \ldots,\ 122,\ 123$?\newline
\textbf{(в)}~Из книги выпала пачка листов. На первой выпавшей странице стоит номер~144, на последней~--- номер~327. Сколько листов выпало?}

\zad{\bigcirc}{\poleD{а}{б}{в}{г}}{Сколько чисел от~1 до~100\newline
\textbf{(а)}~делятся на~3? \qquad \textbf{(б)}~\emph{не} делятся на~3?\newline
\textbf{(в)}~не делятся ни на~3, ни на~5?\newline
\textbf{(г)}~С какой вероятностью наугад выбранное число от~1 до~100 не делится ни на~3, ни на~5?}

\zad{\bigcirc}{\poleA}{В колоде 36~карт четырёх мастей. Для гадания Аза взяла все пики и все картинки (валетов, дам и королей), а остальные карты отложила. Сколько карт выбрала Аза?}

\zad{\bigcirc}{\poleB{а}{б}}{Сколько существует двузначных чисел, у которых\newline
\textbf{(а)}~первая цифра меньше второй? \qquad \textbf{(б)}~сумма цифр двузначна?}

\zadw{\bigcirc}{2.6cm}{Перечислите все тройки \emph{различных} натуральных чисел, дающих в сумме~10. Тройки, отличающиеся только порядком, считаются одинаковыми.}

\vspace{10pt}

\zad{\triangle}{\poleA}{Сколько всего целых чисел от~$m$ до~$n$? (Сами числа $m$ и~$n$ тоже считаются.) Выведите общую формулу.}

\zad{\triangle}{\poleA}{У скольких чисел от~1 до~2019 в записи есть по крайней мере две различные цифры?}

\zad{\triangle}{\poleA}{Среди чисел $1,\ 2,\ \ldots,\ 300$ Валя подчеркнула все те, у которых общий делитель с числом~99 больше~1. Сколько всего чисел она подчеркнула?}

\zadw{\triangle}{2.6cm}{Перечислите все четвёрки \emph{различных} натуральных чисел, дающих в сумме~15. Четвёрки, отличающиеся только порядком, считаются одинаковыми.}

\vtoraya

\zad{}{\poleB{а}{б}}{Сколькими способами можно набрать монетами по 1,~2,~5 и 10~рублей сумму\newline
\textbf{(а)}~10~рублей? \qquad \textbf{(б)}~20~рублей?\newline
Способы, отличающиеся только порядком монет, считаются одинаковыми.}

\zad{}{\poleB{а}{б}}{У каждого из тридцати шестиклассников есть одна ручка, один карандаш и одна линейка. После олимпиады оказалось, что 26~учеников потеряли ручку, 23~--- линейку и 21~--- карандаш.\newline
\textbf{(а)}~Сколько учеников сохранило ручку? линейку? карандаш?\newline
\textbf{(б)}~Какое наименьшее число шестиклассников могло потерять все три предмета?}

\zad{}{\poleC{а}{б}{в}}{Сколько раз встречается цифра~7 в записи чисел\newline
\textbf{(а)}~от~1 до~200? \qquad \textbf{(б)}~от~1 до~2000?\newline
\textbf{(в)}~Сколько \emph{чисел} от~1 до~2000 содержат хотя бы одну семёрку?}

\zad{}{\poleC{а}{б}{в}}{Двенадцать малышей вышли во двор играть в песочнице. Каждый, кто принёс ведёрко, принёс и совочек. Забыли дома ведёрко девять малышей, забыли дома совочек двое.\newline
\textbf{(а)}~Сколько малышей принесли ведёрко? \quad \textbf{(б)}~Сколько принесли совочек?\newline
\textbf{(в)}~На сколько меньше малышей принесли ведёрко, чем совочек?}

\zad{}{\poleB{а}{б}}{Найдите число решений уравнения в целых неотрицательных числах:\newline
\textbf{(а)}~$x+2y+5z=17$; \qquad \textbf{(б)}~$x+2y+5z=37$.}

\zad{}{\poleB{а}{б}}{Каждый из трёх игроков записывает 10~слов, после чего записи сравнивают. Если слово встретилось хотя бы у двоих, его вычёркивают из всех списков.\newline
\textbf{(а)}~Могло ли случиться так, что у первого осталось 8~слов, у второго~--- 7, а у третьего~--- 2?\newline
\textbf{(б)}~А сколько слов \emph{могло} остаться у третьего, если у первого осталось 8, а у второго~--- 7?}

\zad{}{\poleA}{Сумма пяти целых чисел равна~2026. Сколько среди них может быть нечётных? Перечислите все возможности.}

\zad{}{\poleB{а}{б}}{Выбираются пары чисел из набора $1,\ 2,\ \ldots,\ 31$. У скольких пар сумма входящих в них чисел\newline
\textbf{(а)}~нечётна? \qquad \textbf{(б)}~чётна?}

\zad{}{\poleB{а}{б}}{На каждой перемене Робин-Бобин съедает по шоколадке. Уроки были каждый день, кроме воскресений; 1~апреля была среда, а всего за месяц было 100~уроков.\newline
\textbf{(а)}~Сколько было учебных дней? \qquad \textbf{(б)}~Сколько шоколадок он съел?}

\zad{}{\poleB{а}{б}}{Сколько существует \textbf{(а)}~трёхзначных чисел, сумма цифр которых равна~4?\newline
\textbf{(б)}~четырёхзначных чисел, сумма цифр которых больше~33?}

\zadw{}{1.8cm}{Можно ли разложить 99~орехов на 10~кучек так, чтобы любые две кучки отличались, но не более чем в три раза? Если да~--- приведите все такие разложения.}

\zadw{}{1.9cm}{Выпишите все десятизначные числа, у которых первая цифра равна количеству нулей в записи этого числа, вторая~--- количеству единиц, третья~--- количеству двоек, и так далее, а десятая~--- количеству девяток.}

\zadd{}{В классе 25~школьников. Известно, что среди любых троих из них найдутся двое друзей. Докажите, что найдётся школьник, у которого не меньше двенадцати друзей.}

%%%%%%%%%%%%%%%%%%%%%%%%%%%% ВАРИАНТ Б %%%%%%%%%%%%%%%%%%%%%%%%%%%%
\newpage
\shapka{Б}

\zad{\bigcirc}{\poleC{а}{б}{в}}{\textbf{(а)}~Сколько чисел в ряду $23,\ 24,\ 25,\ \ldots,\ 130,\ 131$?\newline
\textbf{(б)}~Сколько чисел в ряду $-23,\ -22,\ \ldots,\ -1,\ 0,\ 1,\ \ldots,\ 130,\ 131$?\newline
\textbf{(в)}~Из журнала выпала пачка листов. На первой выпавшей странице стоит номер~208, на последней~--- номер~451. Сколько листов выпало?}

\zad{\bigcirc}{\poleD{а}{б}{в}{г}}{Сколько чисел от~1 до~150\newline
\textbf{(а)}~делятся на~4? \qquad \textbf{(б)}~\emph{не} делятся на~4?\newline
\textbf{(в)}~не делятся ни на~4, ни на~3?\newline
\textbf{(г)}~С какой вероятностью наугад выбранное число от~1 до~150 не делится ни на~4, ни на~3?}

\zad{\bigcirc}{\poleA}{В колоде 52~карты четырёх мастей. Для фокуса Марк взял все червы и все картинки (валетов, дам и королей), а остальные карты отложил. Сколько карт взял Марк?}

\zad{\bigcirc}{\poleB{а}{б}}{Сколько существует двузначных чисел, у которых\newline
\textbf{(а)}~первая цифра больше второй? \qquad \textbf{(б)}~сумма цифр не больше~5?}

\zadw{\bigcirc}{2.6cm}{Перечислите все тройки \emph{различных} натуральных чисел, дающих в сумме~12. Тройки, отличающиеся только порядком, считаются одинаковыми.}

\vspace{10pt}

\zad{\triangle}{\poleA}{Сколько всего целых чисел от~$a$ до~$b$? (Сами числа $a$ и~$b$ тоже считаются.) Выведите общую формулу.}

\zad{\triangle}{\poleA}{У скольких чисел от~1 до~3000 в записи есть по крайней мере две различные цифры?}

\zad{\triangle}{\poleA}{Среди чисел $1,\ 2,\ \ldots,\ 400$ Гриша подчеркнул все те, у которых общий делитель с числом~65 больше~1. Сколько всего чисел он подчеркнул?}

\zadw{\triangle}{2.6cm}{Перечислите все четвёрки \emph{различных} натуральных чисел, дающих в сумме~16. Четвёрки, отличающиеся только порядком, считаются одинаковыми.}

\vtoraya

\zad{}{\poleB{а}{б}}{Найдите число решений уравнения в целых неотрицательных числах:\newline
\textbf{(а)}~$x+2y+5z=23$; \qquad \textbf{(б)}~$x+2y+5z=43$.}

\zad{}{\poleB{а}{б}}{Выбираются пары чисел из набора $1,\ 2,\ \ldots,\ 41$. У скольких пар сумма входящих в них чисел\newline
\textbf{(а)}~нечётна? \qquad \textbf{(б)}~чётна?}

\zad{}{\poleC{а}{б}{в}}{Пятнадцать ребят вышли на школьный двор. Каждый, кто принёс мяч, принёс и скакалку. Забыли дома мяч одиннадцать ребят, забыли дома скакалку трое.\newline
\textbf{(а)}~Сколько ребят принесли мяч? \quad \textbf{(б)}~Сколько принесли скакалку?\newline
\textbf{(в)}~На сколько меньше ребят принесли мяч, чем скакалку?}

\zad{}{\poleB{а}{б}}{Сколько существует \textbf{(а)}~трёхзначных чисел, сумма цифр которых равна~5?\newline
\textbf{(б)}~четырёхзначных чисел, сумма цифр которых больше~32?}

\zad{}{\poleB{а}{б}}{Сколькими способами можно набрать монетами по 1,~2,~5 и 10~рублей сумму\newline
\textbf{(а)}~15~рублей? \qquad \textbf{(б)}~25~рублей?\newline
Способы, отличающиеся только порядком монет, считаются одинаковыми.}

\zad{}{\poleB{а}{б}}{На каждой перемене Робин-Бобин съедает по шоколадке. Уроки были каждый день, кроме воскресений; 1~марта был понедельник, а всего за месяц было 120~уроков.\newline
\textbf{(а)}~Сколько было учебных дней? \qquad \textbf{(б)}~Сколько шоколадок он съел?}

\zad{}{\poleC{а}{б}{в}}{Сколько раз встречается цифра~4 в записи чисел\newline
\textbf{(а)}~от~1 до~300? \qquad \textbf{(б)}~от~1 до~3000?\newline
\textbf{(в)}~Сколько \emph{чисел} от~1 до~3000 содержат хотя бы одну четвёрку?}

\zad{}{\poleB{а}{б}}{У каждого из двадцати восьми семиклассников есть один шарик, один ластик и одна точилка. После похода оказалось, что 24~ученика потеряли шарик, 21~--- ластик и 18~--- точилку.\newline
\textbf{(а)}~Сколько учеников сохранило шарик? ластик? точилку?\newline
\textbf{(б)}~Какое наименьшее число семиклассников могло потерять все три предмета?}

\zad{}{\poleB{а}{б}}{Каждый из трёх игроков записывает 12~слов, после чего записи сравнивают. Если слово встретилось хотя бы у двоих, его вычёркивают из всех списков.\newline
\textbf{(а)}~Могло ли случиться так, что у первого осталось 9~слов, у второго~--- 8, а у третьего~--- 3?\newline
\textbf{(б)}~А сколько слов \emph{могло} остаться у третьего, если у первого осталось 9, а у второго~--- 8?}

\zad{}{\poleA}{Сумма семи целых чисел равна~2027. Сколько среди них может быть нечётных? Перечислите все возможности.}

\zadw{}{1.8cm}{Можно ли разложить 99~орехов на 10~кучек так, чтобы любые две кучки отличались, но не более чем в три раза? Если да~--- приведите все такие разложения.}

\zadw{}{1.9cm}{Выпишите все десятизначные числа, у которых первая цифра равна количеству нулей в записи этого числа, вторая~--- количеству единиц, третья~--- количеству двоек, и так далее, а десятая~--- количеству девяток.}

\zadd{}{В классе 27~школьников. Известно, что среди любых троих из них найдутся двое друзей. Докажите, что найдётся школьник, у которого не меньше тринадцати друзей.}

\end{document}
